\chapter{Stochastic Cloning and the Sliding Window}
\label{ch:cloning}

\begin{objectives}
\begin{itemize}
  \item See, on a scalar example, why relative measurements require cloned states with correlations.
  \item Derive the general state-augmentation formula $P \gets JPJ\T$.
  \item Derive the project's cloning Jacobian and the invariance property of cloning.
  \item Marginalize the oldest clone in place without allocation, and choose a window policy and size.
\end{itemize}
\end{objectives}

% ------------------------------------------------------------------------------
\section{A scalar example}

Let $x$ be a scalar position that performs a random walk, $x_{k+1} = x_k + w_k$ with $\Var(w_k) = Q$, and suppose its
current estimate has variance $P$. At time 1 we \emph{clone} it, $x_c := x_1$; between times 1 and 2 the state moves; then a
sensor measures the displacement:
\begin{equation}
  z = x_2 - x_c + n, \qquad \Var(n) = R .
\end{equation}
Augmenting the state with the clone gives the vector $(x, x_c)$. At the moment of cloning both entries are the \emph{same}
random variable, so their joint covariance is $\begin{bmatrix} P & P \\ P & P\end{bmatrix}$, which is singular. Propagating $x$ but
not $x_c$ gives
\begin{equation}
  P_{\mathrm{aug}} = \begin{bmatrix} P + Q & P \\ P & P \end{bmatrix}, \qquad H = \begin{bmatrix} 1 & -1 \end{bmatrix}.
\end{equation}
The innovation covariance and gain are
\begin{equation}
  S = (P + Q) - 2P + P + R = Q + R, \qquad K = \frac{1}{Q + R}\begin{bmatrix} Q \\ 0 \end{bmatrix}.
\end{equation}
The result is exactly right: a displacement measurement is informative only about what happened \emph{between} the two
times, so $S$ involves only $Q$ and $R$. The update corrects the current state by the part of its uncertainty that arose
after the clone, and leaves the clone untouched. The posterior variance of $x_2$ is $P + Q - Q^2/(Q + R)$, which stays
large if $P$ is large, as it must: a relative measurement cannot reveal absolute position.

Now suppose we had treated the past estimate as known (no clone, zero variance). Then $S = P + Q + R$ and the posterior
variance of $x_2$ would be $(P + Q)R/(P + Q + R)$, which tends to $R$ when $P$ is large. The filter would believe it had
learned the absolute position to within the sensor noise from a displacement measurement: a severe overconfidence.
Treating the past estimate as an independent random variable is wrong in a different way. Only the exact
cross-covariance gives the right answer, and stochastic cloning~\cite{roumeliotis2002cloning} is the mechanism that keeps it.

% ------------------------------------------------------------------------------
\section{State augmentation}

\begin{proposition}[Augmentation with a function of the state]
\label{prop:augmentation}
Let $x$ have error covariance $P$, and append a new variable $y = g(x)$ whose error satisfies $\delta y \approx J_g\,\delta x$. The
augmented error $(\delta x, \delta y) = J\,\delta x$ with $J = \begin{bmatrix} I \\ J_g\end{bmatrix}$ has covariance
\begin{equation}
  P \gets J P J\T = \begin{bmatrix} P & PJ_g\T \\ J_gP & J_gPJ_g\T \end{bmatrix}.
  \label{eq:augmentation}
\end{equation}
\end{proposition}
\begin{proof}
\Cref{prop:affine} with the linear map $J$.
\end{proof}

The augmented covariance has rank $\dim x$: it is singular, because $y$ is a deterministic function of $x$ at the moment of
augmentation. This is correct and harmless. Propagation adds process noise to the IMU block only, after which the clone and
the current state are no longer perfectly correlated. The Joseph form and symmetrization keep the matrix positive
semi-definite throughout.

\section{The project's clone}

At an image timestamp, after propagating the IMU state to exactly that time, the filter appends a clone holding the current
pose $(p_{wb}, q_{wb})$, and augments the covariance with the cloning Jacobian (AGENTS.md F5).

\begin{derivex}{the cloning Jacobian $J$}{sol:clone}
The clone's error is defined with the same retraction as the IMU pose (AGENTS.md C4): $p_c = \hat p_c + \delta p_c$ and
$q_c = \hat q_c\qprod\Exp_q(\delta\theta_c)$, with $\hat p_c = \hat p$ and $\hat q_c = \hat q$ at cloning time.
\begin{enumerate}
  \item At the moment of cloning the true clone equals the true IMU pose. Express $\delta p_c$ and $\delta\theta_c$ in terms of the IMU
        error state. Is any linearization needed?
  \item Write $J_g \in \R^{6\times(15+6N)}$ and then $J \in \R^{(21+6N)\times(15+6N)}$, with the new clone appended at the end of the
        state.
  \item What would change if the clone stored the \emph{camera} pose $(p_{wc}, q_{wc})$ instead of the IMU pose? (The project
        deliberately does not do this.)
\end{enumerate}
\textbf{Check:} Milestone 10a requires that after cloning, the new clone's covariance block and its cross-covariances equal the
corresponding IMU pose blocks \emph{exactly}, and that a run with cloning enabled but only pose updates reproduces the
IMU estimate and covariance of a run without cloning to $10^{-10}$.
\end{derivex}

\begin{proposition}[Cloning does not change the IMU marginal]
\label{prop:clone-invariance}
Augmentation, propagation with static clones, and marginalization of clones leave the marginal distribution of the IMU state
unchanged, provided no measurement involves the clones.
\end{proposition}
\begin{proof}
Augmentation leaves the $P_{II}$ block unchanged by \eqref{eq:augmentation}. Propagation updates $P_{II}$ by $\Phi P_{II}\Phi\T + Q_d$
regardless of the clones. Marginalization deletes clone rows and columns. A measurement that involves only the IMU state
has a Jacobian with zero clone columns; its gain for the IMU block, $P_{II}H_I\T S\inv$, and $S = H_IP_{II}H_I\T + R_z$ do not
involve the clones either.
\end{proof}

This is the invariance test of Milestone 10a. It is a powerful end-to-end check of the augmentation, propagation and
marginalization code, because any indexing error in the clone blocks leaks into the IMU block and breaks the equality.

% ------------------------------------------------------------------------------
\section{Clones in propagation and updates}

Clones are static (\cref{ch:errordyn}): propagation multiplies their cross-covariance with the IMU state by $\Phi$ and leaves the
clone-clone block alone. Visual measurements depend only on clone states, so their Jacobians are zero in the IMU columns.
The update nevertheless corrects the IMU state, through the gain rows $P_{IC}H_C\T S\inv$: information about relative motion
between past poses corrects the current velocity, attitude and biases through the correlations built up by propagation.
Injection and reset (\cref{sec:reset}) apply to every clone, each with its own $I - \tfrac12\sx{\delta\hat\theta_i}$ block.

% ------------------------------------------------------------------------------
\section{Marginalizing the oldest clone}

When the window holds \code{kMaxClones} clones and a new image arrives, the oldest clone must go. In covariance form,
marginalization is deletion of its six rows and columns (\cref{thm:gauss-cond}). Before deleting it, the filter must process
every feature that has an observation at that clone (AGENTS.md F5), because those observations become unusable once
the clone is gone.

\begin{implnote}[in-place deletion without allocation]
$P$ is stored in a fixed-max-size matrix whose active part is the top-left $n\times n$ corner. Deleting the block at offset $o$
means shifting every row below it up by six and every column to its right left by six, then reducing $n$ by six. The source
and destination regions overlap. Copying toward lower indices in ascending order is safe, because each element is read
before it is overwritten (\code{std::copy} permits overlapping ranges when the destination begins before the source), but an
Eigen block assignment between overlapping blocks of the same matrix is an aliasing hazard whose traversal order is not
guaranteed. Write the shift explicitly, or through a preallocated member workspace, and test it against a reference built
with a dynamic matrix (Milestone 10a).
\end{implnote}

The clone states themselves are stored in a fixed-capacity container, such as a ring buffer of $N_{max}$ poses with
timestamps, so that removing the oldest is a pointer move. If the covariance blocks are kept in the same order as the
container, a ring-buffer layout could even avoid the shift, at the cost of more complex index arithmetic. Simplicity is
usually the better choice here: the shift costs $\mathcal{O}(n^2)$ once per image.

% ------------------------------------------------------------------------------
\section{Window policy and size}

The project uses the simplest policy: first in, first out. Each image creates a clone, and the oldest is marginalized when
the window is full. More elaborate policies keep clones that are geometrically useful: if the vehicle barely moved since the
previous image, marginalize the second-newest clone instead of the oldest, so that the window spans a longer baseline.

The window size trades accuracy for cost.
\begin{itemize}
  \item \textbf{Longer windows} let longer tracks be used in full, with more parallax, and give each update more
        constraints.
  \item \textbf{Cost} grows as $\mathcal{O}(n^3)$ per update with $n = 15 + 6N$, and as $\mathcal{O}(n^2)$ per IMU step for the
        cross-covariance propagation.
  \item \textbf{Memory}: a fixed-max covariance of doubles occupies $8(15 + 6N_{max})^2$ bytes: 45~kB for $N_{max} = 10$, 146~kB for
        20, 304~kB for 30. Workspaces for $H$, $S$ and $K$ add similar amounts.
\end{itemize}

\begin{pitfall}[the stack budget]
Fixed-max-size Eigen matrices store their maximum-size buffer \emph{inline}, with no heap allocation. That is what makes them
allocation-free, but it also means a filter object constructed on the stack carries all of its matrices on the stack. With
the default 1~MB main-thread stack on Windows, a filter with a few hundred kilobytes of workspaces can overflow it. Construct
the filter once, during the load phase, as a heap object (for example with \code{std::make\_unique}), which the memory tiers
allow, and keep the large matrices as members rather than function-local temporaries (AGENTS.md E4, Milestone 10a).
\end{pitfall}

\begin{takeaways}
\begin{itemize}
  \item A relative measurement between a past and a present state needs the past state in the filter, with its exact
        correlations. Treating the past estimate as known makes the filter overconfident.
  \item Augmentation with $y = g(x)$ uses $P \gets JPJ\T$, $J = [I; J_g]$; the result is singular at first and becomes regular
        after propagation.
  \item Cloning, static propagation and marginalization leave the IMU marginal unchanged; this is a strong test.
  \item Marginalizing a clone deletes its rows and columns, after first processing its features. Shift in place carefully.
  \item Fixed-max matrices live inline: keep them as members of a heap-allocated filter.
\end{itemize}
\end{takeaways}

\begin{exercises}
  \item In the scalar example, compute the posterior covariance of $(x_2, x_c)$ and verify that the clone's variance is
        unchanged. When would a relative measurement change the clone's estimate?
  \item Extend the scalar example to two clones and two displacement measurements, and show that processing them gives the
        same result as one measurement of $x_2 - x_{c_1}$ followed by one of $x_{c_2} - x_{c_1}$, if correlations are kept.
  \item Implement in-place deletion of a $6\times6$ block from a $75\times75$ symmetric matrix stored in a larger buffer, and test it
        against a reference.
  \item For $N_{max} = 10, 20, 30$, estimate the flop count of one MSCKF update with 50 features of 6 observations each, with and
        without QR compression (\cref{ch:msckfupdate}).
\end{exercises}
